% Section 12.4, from lectures/L12/appendix/c_exercises.md.

\section{Exercises}
\label{c12:sec:exercises}

\begin{aside}
Exercise~\ref{c12:ex:metaprev} consolidates \secref{c12:sec:metaprev};
Exercise~\ref{c12:ex:bittimer} to \ref{c12:ex:drift} consolidate \secref{c12:sec:bittimer}. Design
each module from the specification in its own section; the testbenches handed out, run per
appendix~\ref{app:sim}, are the check that they behave correctly.
\end{aside}

% ------------------------------------------------------------------------------------------
\exerciseset{The synchroniser}

\exercise{\code{meta_prev}}{VHDL}
\label{c12:ex:metaprev}
You have built this module before, so most of what is here is about the two things that are new.

\xpart{a} Write your own \file{meta_prev.vhd} from \secref{c12:sec:metaprev} and run its testbench:

\begin{shell}
cd controller
ghdl -a --std=93 meta_prev.vhd meta_prev_tb.vhd
ghdl -e --std=93 meta_prev_tb
ghdl -r --std=93 meta_prev_tb --assert-level=error
\end{shell}

Note that no \file{can_def.vhd} appears on the first line. \code{meta_prev} is the only module in the
book that reads no package at all, which is also why it is the one reused three times.

\xpart{b} Remove one of the two flip-flops, leaving a single register, and rerun. Which check fails,
and what does its message tell you? Put it back. (A one-flip-flop version is not merely weaker; it is
exactly the bug the testbench exists to catch, and it would otherwise pass every functional
simulation in the book, since nothing in a simulator is ever metastable.)

\xpart{c} \code{meta_prev} has no reset. Say what \code{sync_out} reads at the first clock edge after
power-up, why that is harmless here, and what would go wrong in \code{can_controller} itself if the
module \emph{had} had a \code{reset_s2_n} input: the \code{reset_sync} instance is the one producing
\code{reset_s2_n}, so what would it take for its own reset?

\xpart{d} \code{can_controller} instantiates it twice, both at the default width. Say which two of
its fifteen ports that concerns, and why those two and no others. Then say what the design's third
instance synchronises \textemdash{} the one \code{can_spi_node} owns at its top level in
\chapref{c19:ch} \textemdash{} and why that one in particular cannot sit inside
\code{can_controller} instead.

\xpart{e} Two switches are flipped at the same instant into an instance with \code{WIDTH = 10}.
Explain why \code{sync_out} can briefly show one switch's new value and the other's old one, why that
is not a bug in \code{meta_prev}, and why it would matter considerably more if the ten bits were a
counter being read across a clock domain boundary instead of switches set by hand.

% ------------------------------------------------------------------------------------------
\exerciseset{The bit timer}

\exercise{\code{bit_timer}}{VHDL}
\label{c12:ex:bittimer}
Write your own \file{bit_timer.vhd} from \secref{c12:sec:bittimer}. Then verify it:

\begin{shell}
cd controller
ghdl -a --std=93 can_def.vhd bit_timer.vhd bit_timer_tb.vhd
ghdl -e --std=93 bit_timer_tb
ghdl -r --std=93 bit_timer_tb --assert-level=error
\end{shell}

If a check fails, read the assertion message (appendix~\ref{app:sim}); it names exactly which tick
and which expected value it concerns.

\exercise{A configurable resynchronisation offset}{Simulation}
\label{c12:ex:offset}
Real CAN hardware sometimes needs a bit timer that can resynchronise \emph{and} immediately jump
forward a fixed number of ticks (a simplified stand-in for the phase segment adjustments real CAN bit
timing logic performs, mentioned in \secref{c10:sec:bittiming}).

\xpart{a} Add a new generic, \code{RESYNC_OFFSET: natural := 0}, to \code{bit_timer}. When
\code{resync = '1'} the counter should be set to \code{RESYNC_OFFSET} (instead of always \code{0}) at
the next clock edge.

\xpart{b} With \code{RESYNC_OFFSET = 0} your modified entity has to behave exactly like the original;
confirm that by running the \emph{unchanged} \file{bit_timer_tb.vhd} against it; it should still pass
with no changes to the testbench.

\xpart{c} Work out on paper where \textbf{both} outputs land when \code{RESYNC_OFFSET} is \code{10}:
how many ticks after a \code{resync} pulse \code{sample} fires, and how many ticks after it
\code{bit_done} does. Both move, and for the same reason \textemdash{} the counter restarts at
\code{10} instead of \code{0}, so the whole period is ten ticks shorter.

Then let the testbench handed out measure it for you. Temporarily change your generic's default value
to \code{10} and rerun \code{bit_timer_tb}. It binds no generic map, so it picks up whatever default
you set, and case 4 is the one that notices:

\begin{console}
bit_timer_tb.vhd:132:17:@6832ns:(assertion error): bit_timer: unexpected sample at tick 25 after resync!
\end{console}

Note \textbf{which} of the two checks fails first. Case 4 checks not only that \code{bit_done} comes
at the right tick but also that the sample point moves with it, and the sample point comes first in
the period \textemdash{} so it is the one that fires, fourteen ticks before the \code{bit_done} check
would have. Under \code{--assert-level=error} the simulation stops there, so you never see the
message about \code{bit_done}.

The message names the tick itself: \code{sample} came at tick 25, that is,
\code{SAMPLE_TICK - RESYNC_OFFSET}. If you want to see \code{bit_done} too, lower the sample check's
\code{severity error} to \code{severity note} in your own copy of the testbench and rerun; it then
reports tick 39, that is, \code{TICKS_PER_BIT - 1 - RESYNC_OFFSET}. Confirm that both agree with your
paper answer, then put the default back to \code{0} and check that the \emph{unchanged} testbench
passes again.

That is worth noticing in itself: a testbench that fails \emph{usefully} says not only that a design
is wrong but by how much, which is why the assertion messages in this project all name the tick or
the value they expected.

\exercise{How long can a bit timer free-run?}{Reasoning}
\label{c12:ex:drift}
\code{bit_timer} resynchronises \textbf{once per frame}, at SOF, and free-runs from there
(\secref{c10:sec:stuffing}'s clock recovery argument). That works only if the node's own oscillator
stays close enough to every other node's for the whole frame. This exercise puts a number on ``close
enough''.

Suppose a node's 50 MHz oscillator in fact runs 1~\% fast, so that its ticks are 1~\% shorter than
the transmitter's, and that everything else is as specified: \code{TICKS_PER_BIT = 50},
\code{SAMPLE_TICK = 35}.

\xpart{a} How many ticks has this node's counter gained on the transmitter's after one bit period?
(Express it as a fraction of one tick.)

\xpart{b} The node samples at tick 35 of 50, so relative to the true bit it can drift 15 ticks
\emph{later} before it samples past the bit's end, or 35 ticks \emph{earlier} before it samples
before the bit begins. A fast clock drifts the sample point earlier. After how many bit periods has
it drifted the full 35 ticks?

\xpart{c} A standard frame with 8 data bytes runs to about 110--130 bits on the wire once the stuff
bits are counted. Compare that with your answer to \textbf{b)}: is this node still sampling correctly
at EOF? At roughly which field does it start sampling the wrong bit?

\xpart{d} Real CAN controllers resynchronise \emph{during} a frame, not only at SOF, using the phase
segments and the synchronisation jump width \secref{c10:sec:bittiming} mentions. Explain in one or
two sentences what that buys them, expressed in terms of your answer to \textbf{b)}.
